\begingroup
\StyledTableSetup
\scriptsize
\setlength{\tabcolsep}{3pt}
\newcolumntype{L}[1]{>{\RaggedRight\arraybackslash\hyphenpenalty=10000\exhyphenpenalty=10000\emergencystretch=1em}p{#1}}
\newcolumntype{Z}{>{\RaggedRight\arraybackslash\hyphenpenalty=10000\exhyphenpenalty=10000\emergencystretch=1em}X}
\newcommand{\titreTableauCrisesMajeures}{Évolution des vulnérabilités et des crises majeures contemporaines en France : 1982 à 2025.}
\newcommand{\noteCrisesTerrorismeAvantDeuxMilleTexte}{Voir notamment l'intervention de Monsieur Jean-Jacques URVOAS garde des Sceaux, ministre de la Justice, le $\text{1}^{\text{er}}$ juin 2016 : \cite{renforcement_lutte_terrorisme_82_86}.}

\label{tab:crises_majeures_page_debut}%
\begin{xltabular}{\linewidth}{@{}
L{1.1cm}
L{2.35cm}
L{2.25cm}
Z
Z
Z
Z
@{}}
\toprule
\textbf{Période} &
\textbf{Crise emblématique} &
\textbf{Nature dominante} &
\textbf{Vulnérabilités révélées} &
\textbf{Effets sur la société et la cohésion nationale} &
\textbf{Évolutions des réponses publiques} &
\textbf{Enjeux de résilience et rôle potentiel des citoyens engagés} \\
\midrule
\endfirsthead

\toprule
\textbf{Période} &
\textbf{Crise emblématique} &
\textbf{Nature dominante} &
\textbf{Vulnérabilités révélées} &
\textbf{Effets sur la société et la cohésion nationale} &
\textbf{Évolutions des réponses publiques} &
\textbf{Enjeux de résilience et rôle potentiel des citoyens engagés} \\
\midrule
\endhead

\midrule
\multicolumn{7}{r}{\small\itshape Suite page suivante} \\
\endfoot

\endlastfoot

1982 - 1986 &
Vague d'attentats à Paris &
Terroriste &
Vulnérabilité des espaces urbains et civils &
Climat d'insécurité et de peur au sein de la population &
Renforcement des dispositifs antiterroristes\footnote{\noteCrisesTerrorismeAvantDeuxMilleTexte}\xdef\noteCrisesTerrorismeAvantDeuxMille{\number\value{footnote}} &
Sensibilisation citoyenne et vigilance collective \\
\addlinespace

1995 &
Attentats du RER Saint-Michel &
Terrorisme intérieur &
Fragilité des transports publics et nécessité du renseignement &
Traumatisme collectif durable et inquiétude nationale &
Modernisation de la coordination sécuritaire\footnotemark[\noteCrisesTerrorismeAvantDeuxMille] &
Importance de la diffusion d'informations fiables \\
\addlinespace

1999 &
Tempêtes Lothar et Martin &
Catastrophe naturelle &
Dépendance énergétique et fragilité des réseaux critiques &
Désorganisation territoriale massive et interruption des services essentiels &
Renforcement des infrastructures et plans de prévention\footnote{Voir notamment le rapport d'étape de la mission interministérielle sur l'évaluation des dispositifs de secours et d'intervention suite aux tempêtes du 26 et 28 décembre 1999 : \cite{renforcement_dispositif_prevention_1999}.} &
Développement d'une culture locale de préparation aux crises \\
\addlinespace

2003 &
Canicule &
Sanitaire et climatique &
Isolement des personnes vulnérables et saturation hospitalière &
Crise sanitaire majeure et émotion nationale &
Création du plan national canicule\footnote{Voir notamment l'instruction ministérielle du 27 mai 2004 relative à la gestion sanitaire des vagues de chaleur en France métropolitaine : \cite{gestion_sanitaire_des_vagues_de_chaleur_2004}.} &
Mobilisation des réseaux associatifs et solidarité de proximité \\
\addlinespace

2005 &
Émeutes urbaines &
Sociale et sécuritaire &
Fragilités territoriales et tensions sociales persistantes &
Violences urbaines et défiance envers les institutions &
Renforcement des politiques de cohésion sociale et de sécurité\footnote{Accélération des plans banlieues dont le premier fut initié en 1973 : \cite{Le_Monde_plans_banlieue}.} &
Rôle de médiation et de prévention locale \\
\addlinespace

2010 &
Tempête Xynthia &
Catastrophe naturelle &
Exposition des territoires littoraux et urbanisation vulnérable &
Pertes humaines et destructions importantes &
Révision des politiques de prévention des submersions marines &
Sensibilisation des collectivités et préparation des populations\footnote{\cite{Senat:2015:Xynthia}.} \\
\addlinespace

2015 &
Attentats terroristes de Paris &
Terrorisme hybride &
Vulnérabilité sécuritaire et informationnelle &
Traumatisme national et instauration de l'état d'urgence &
Renforcement de Vigipirate et de la législation antiterroriste &
Résilience psychologique collective et lutte contre les infox\footnote{\cite{AN:2016:Lutte_terrorisme}.} \\
\addlinespace

2015 - 2016 &
Crise migratoire de Calais &
Humanitaire et géopolitique &
Saturation des capacités d'accueil et tensions territoriales &
Polarisation du débat public et tensions locales &
Réorganisation des dispositifs d'accueil et de contrôle &
Importance du tissu associatif et du dialogue territorial\footnote{\cite{AN:2021:Migration}.} \\
\addlinespace

2018 - 2019 &
Crise des Gilets jaunes &
Sociale et sociétale &
Fragilité du lien entre institutions et population &
Violences, polarisation et défiance institutionnelle &
Grand débat national et mesures socio-économiques &
Nécessité d'une médiation territoriale et citoyenne\footnote{\cite{AN:2019:Gilets_Jaunes}.} \\
\addlinespace

2020 &
Pandémie de Covid-19 &
Sanitaire systémique &
Dépendances logistiques, hospitalières et économiques &
Confinements, isolement, défiance et crise économique &
Coordination interministérielle et plans de relance &
Rôle essentiel des réseaux citoyens, associatifs et informationnels\footnote{\cite{AN:2020:Covid}.} \\
\addlinespace

2022 &
Guerre en Ukraine et effets en France &
Géopolitique et hybride &
Dépendances énergétiques, cyber et informationnelles &
Inflation, tensions énergétiques et inquiétudes stratégiques &
Renforcement de la résilience énergétique et cybersécuritaire &
Développement de la culture de défense et de résilience nationale\footnote{\cite{economie_gouv_resilience_ukraine}.} \\
\addlinespace

2023 &
Réforme des retraites &
Sociale et politique &
Tensions autour de la gouvernance et de la représentation &
Manifestations et crispations sociales persistantes &
Usage du 49.3 et tensions démocratiques accrues &
Importance du dialogue civique et de l'acceptabilité sociale \\
\addlinespace

2023 &
Émeutes après la mort de Nahel &
Sociale, sécuritaire et informationnelle &
Fragilité du lien police-population et amplification via les réseaux sociaux &
Violences urbaines rapides et diffusion massive d'images &
Réflexion sur la police de proximité et l'autorité publique &
Besoin d'acteurs de médiation et de relais locaux crédibles \\
\addlinespace

\pagebreak
2024 - 2025 &
Crises climatiques : sécheresses et incendies &
Environnementale systémique &
Pression sur les ressources, l'agriculture et les territoires &
Restrictions, tensions économiques et environnementales &
Plans d'adaptation climatique et politiques de résilience &
Développement d'une culture durable de préparation territoriale \\
\bottomrule

\end{xltabular}
\par\nopagebreak\vspace{0.4em}
\begin{center}
\normalsize
\refstepcounter{table}%
\label{tab:crises_majeures}%
\label{tab:crises_majeures_page_fin}%
\addtocontents{lot}{\protect\contentsline{table}{\protect\numberline{\thetable}\titreTableauCrisesMajeures}{\protect\pageref{tab:crises_majeures_page_debut}--\protect\pageref{tab:crises_majeures_page_fin}}{table.\thetable}}%
Tableau~\thetable{} : \titreTableauCrisesMajeures
\end{center}
\endgroup
